\chapter{Conclusiones}
\label{sec:conclusiones}

El presente trabajo abordó el diseño analítico, la simulación circuital y la
validación experimental en laboratorio de un sistema de acondicionamiento para
señales electrocardiográficas (ECG). La arquitectura desarrollada integró un
amplificador de instrumentación de tres amplificadores operacionales con
filtrado pasabanda activo en su propia etapa restadora, complementado con una
etapa final de post-amplificación y calibración. En lo relativo a la ganancia
diferencial, la implementación cumplió holgadamente la cota impuesta ($\ge
\SI{60}{\decibel}$): 
se registró una ganancia de $\SI{66.77}{\decibel}$ ($2180.0\ \text{V/V}$) a
$\SI{25}{\hertz}$. Esta amplificación transforma señales cardíacas débiles
($\sim \SI{1}{\milli\volt\text{pp}}$ a $\SI{2.5}{\milli\volt\text{pp}}$) en
amplitudes de $\SI{2}{\volt}$ a $\SI{3.5}{\volt\text{pp}}$, ideales
para su digitalización sin distorsión armónica apreciable.

La incorporación de reactancias en la etapa restadora resultó altamente eficaz
para la selectividad espectral del circuito, suprimiendo la necesidad de etapas
de filtrado adicionales. La función de transferencia presentó un cero en el
origen que garantizó ganancia nula a corriente continua ($A_v(0) = 0$),
confiriendo inmunidad absoluta frente a potenciales galvánicos de contacto
electrodo-piel (hasta $\pm \SI{300}{\milli\volt}$) y suprimiendo cualquier
corrimiento de línea de base. Asimismo, el polo inferior en $f_L \approx
\SI{0.07}{\hertz}$ y el corte superior nominal en $f_{H,\text{sim}} \approx
\SI{72.3}{\hertz}$ se contrastaron con la placa física, donde se midió un corte
experimental de $f_{H,\text{med}} = \mathbf{\SI{82.2}{\hertz}}$ y un ancho de
banda efectivo de $\mathbf{\SI{82.2}{\hertz}}$, plenamente apto para el
monitoreo ECG. 

En cuanto a la inmunidad frente a ruidos de modo común, el prototipo físico
alcanzó un $\text{RRMC}_{\text{med}} = \mathbf{\SI{98.80}{\decibel}}$ a
$\SI{25}{\hertz}$ ($A_{md} = 2180.0\ \text{V/V}$ y $A_{mc} = 0.0250\
\text{V/V}$), superando con holgura la especificación exigida ($\gg
\SI{80}{\decibel}$). Este resultado asegura una eliminación contundente de las
interferencias de red de $\SI{50}{\hertz}$. Por último, la elección de
amplificadores BiFET (TL084) fue plenamente acertada: su elevada impedancia de
entrada  y sus corrientes de polarización mínimas, sumando su disponibilidad en
tiendas locales de electrónica y su buen costo económico contribuyeron a la
satifactoria realización del trabajo práctico. En conjunto, la estrecha
correlación entre cálculo teórico, simulación y mediciones sobre PCB valida la
robustez, compacidad y viabilidad del diseño propuesto.
